% =====================================================================
%  TP4, Phase III, supplementary material.
%  Author: FOMETHE SOBMBANANG MAXIMILIEN
%
%  OPTIONAL. The eight page article allocates one page to the protocol and
%  two to the results; this file holds the material that does not fit in
%  that budget but that supports claims made in the main text. Include it
%  only if space allows, or keep it as supplementary material. Dropping it
%  leaves the main text self-contained: the claims it supports are stated
%  there with their numbers.
% =====================================================================

\appendix

\section{Supplementary material}
\label{sec:appendix}

\subsection{Default cut-off against a cut-off chosen on validation}
\label{sec:appendix-operating}

\Cref{tab:operating} reads the seven baselines at two operating points: the
default cut-off of each family, and the threshold that maximises F1 on
validation and is then transferred unchanged to test. Every model improves and
the degenerate cases disappear. Boosting goes from \num{0.0870} to \num{0.3639},
the RBF machine from \num{0.0000} to \num{0.3411} and the linear machine from
\num{0.0610} to \num{0.3070}, while Gaussian naive Bayes, first at the default
cut-off with \num{0.2995}, moves only to \num{0.3099} and falls back to the
middle of the table. The ordering by F1 at the selected threshold then broadly
follows the ordering by average precision, which is the expected behaviour: once
the threshold stops being an uncontrolled variable, the thresholded metric
measures the same underlying ranking quality.

\begin{table}[t]
  \centering
  \caption{Test precision, recall and F1 at the default cut-off of each family
  (\num{0.5} on a posterior, \num{0} on a margin) and at the cut-off that
  maximises validation F1. The test split is never used to choose the threshold.
  The two models that made almost no positive call at the default cut-off gain
  the most.}
  \label{tab:operating}
  \footnotesize\setlength{\tabcolsep}{4pt}
  \begin{tabular}{lcccccc}
\toprule
\multirow{2}{*}{Model} & \multicolumn{3}{c}{Default cut-off} & \multicolumn{3}{c}{Cut-off chosen on validation} \\
\cmidrule(lr){2-4}\cmidrule(l){5-7}
 & Prec. & Rec. & F1 & Prec. & Rec. & F1 \\
\midrule
Histogram gradient boosting & 0.388 & 0.049 & 0.087 & 0.264 & 0.587 & 0.364 \\
Kernel SVM (RBF) & 0.000 & 0.000 & 0.000 & 0.264 & 0.483 & 0.341 \\
Linear SVM (soft margin) & 0.317 & 0.034 & 0.061 & 0.235 & 0.442 & 0.307 \\
K-nearest neighbours & 0.324 & 0.024 & 0.045 & 0.213 & 0.570 & 0.310 \\
Decision tree (CART) & 0.244 & 0.016 & 0.029 & 0.231 & 0.490 & 0.314 \\
Gaussian Naive Bayes & 0.192 & 0.676 & 0.299 & 0.207 & 0.614 & 0.310 \\
Logistic regression (SGA) & 0.250 & 0.102 & 0.145 & 0.205 & 0.484 & 0.288 \\
\bottomrule
\end{tabular}

\end{table}

\begin{figure}[t]
  \centering
  \includegraphics[width=\columnwidth]{fig5_threshold_sensitivity.pdf}
  \caption{F1 of the boosted model against its decision threshold, on validation
  and on test. The default \num{0.5} cut-off lies far to the right of the
  maximum, which on this prevalence sits near \num{0.07}. The two curves peak at
  nearly the same place, so the threshold transfers across splits.}
  \label{fig:threshold}
\end{figure}

\Cref{fig:threshold} shows why. For the boosted model the F1 maximum lies at a
posterior of \num{0.0691}, not at \num{0.5}, and the validation and test curves
peak at nearly the same place, so the choice transfers. A thresholded metric on
an imbalanced task therefore reports the sum of two things, the quality of the
ranking and the accident of where the library put the cut-off, and on this
prevalence the second term dominates.

\subsection{Does reducing the dimension help?}
\label{sec:appendix-representation}

The two families whose behaviour depends most directly on the ambient dimension
were rerun on a \num{50} component principal projection fitted on the training
subsample, retaining \num{95.3} percent of the variance
(\cref{tab:representation}). Neither improves: the $K=5$ neighbours rule falls
from \num{0.1822} to \num{0.1776} validation average precision and the RBF
machine from \num{0.2417} to \num{0.2344}. The expectation behind the ablation,
that Euclidean distances and an isotropic kernel would behave better once
\num{784} correlated pixel coordinates are compressed, is not met. Part of the
discriminative signal lies in directions that principal component analysis
discards for carrying little variance, which is plausible on this modality:
variance in a chest radiograph is dominated by patient size, framing and
exposure, none of which is the costophrenic angle.

\begin{table}[t]
  \centering
  \caption{Representation check on the primary seed, at matched
  hyperparameters. The principal basis is fitted on the training subsample only
  and retains \num{95.3} percent of the variance.}
  \label{tab:representation}
  \footnotesize\setlength{\tabcolsep}{5pt}
  \begin{tabular}{lcccc}
\toprule
\multirow{2}{*}{Model} & \multicolumn{2}{c}{Pixels (784)} & \multicolumn{2}{c}{PCA (50)} \\
\cmidrule(lr){2-3}\cmidrule(l){4-5}
 & AUC & AP & AUC & AP \\
\midrule
KNN, $K = 5$ & 0.6336 & 0.1819 & 0.6317 & 0.1810 \\
Kernel SVM, $C = 1$ & 0.7175 & 0.2763 & 0.7086 & 0.2683 \\
\bottomrule
\end{tabular}

\end{table}

The projection does buy computation: the RBF machine then scores the test split
in \SI{21.9}{\second} instead of \SI{146.7}{\second}, a factor of \num{6.7}, for
\num{0.008} of average precision. That trade is worth taking in a deployment
setting and is not worth taking when the point is to compare model families, so
the main tables keep the \num{784} dimensional features.

\subsection{Reproducing the numbers}
\label{sec:appendix-repro}

The code, the result files and the MLflow store live in the branch
\texttt{phase3-fomethe} of the project repository. After placing the six
ChestMNIST arrays under \texttt{data/}:

\begin{center}
\footnotesize
\begin{tabular}{@{}ll@{}}
\toprule
Command & Produces \\
\midrule
\texttt{python tests/test\_phase3.py}  & \num{12} invariant checks \\
\texttt{python src/run\_baselines.py}  & \cref{tab:baselines}, \cref{tab:operating} \\
\texttt{python src/run\_robustness.py} & \cref{tab:robustness}, \cref{tab:degradation} \\
\texttt{python src/make\_figures.py}   & \cref{fig:pr} to \cref{fig:threshold} \\
\texttt{python src/export\_tables.py}  & the LaTeX bodies of every table \\
\bottomrule
\end{tabular}
\end{center}

\noindent
No number in the text or in the tables is retyped by hand: the table bodies are
generated from the result files, and the figures are drawn from the score
vectors those same runs persisted, so a figure cannot disagree with the table
beside it.
